\clearpage\begingroup\SourceLeaf{78b}\HistoricalHead{73}{R}{GAUSS AN OLBERS.}\HistoricalFoot{XI1.}{10}
\begin{center}[3.]\end{center}
\begin{center}\textsc{Gauss} an \textsc{Schumacher}.\quad Göttingen, den 3. März 1832.\end{center}
\begin{center}
\rule{.70\textwidth}{.4pt}\\
{\small [Briefwechsel zwischen \textsc{Gauss} und \textsc{Schumacher} II, S. 295.]}\\[-.3ex]
\rule{.70\textwidth}{.4pt}
\end{center}
\noindent . . . . . . . . . . . . . . . . . . . . .\par
\indent Ich bin, wie Sie leicht denken können, zu wissenschaftlichen Arbeiten lange
Zeit wenig aufgelegt gewesen, habe aber doch in der letzten Zeit ein ziemlich
lebhaftes Interesse für einen Gegenstand gewonnen, oder vielmehr erneuert,
denn von jeher habe ich denselben als einen sehr reichhaltigen betrachtet,
aber erst jetzt ist mir alles, was mir früher darin dunkel war, in grosse Klarheit
getreten. Dies ist der Erdmagnetismus, und ich möchte wohl Ihre Verwendung
ansprechen, um einen Wunsch in Erfüllung gehen zu sehen. Der
vortreffliche \textsc{Hansteen}[\textsuperscript{*})]\begingroup\renewcommand{\thefootnote}{}
\footnotetext{[\textsuperscript{*}) \textsc{Chr. Hansteen}, \textit{Isodynamische Linien für die ganze Magnetkraft der Erde},
\textsc{Poggendorffs} Annalen der Physik u. Chemie \textbf{9}, 1827, S. 49.]}
\endgroup\ hat uns vor einiger Zeit eine Karte der isodynamischen
Linien geliefert, und hoffentlich haben wir von demselben auch bald neue
\clearpage\SourceLeaf{79}\HistoricalHead{74}{L}{PHYSIK. BRIEFWECHSEL.}
\noindent Declinations- und Inclinationskarten zu erwarten. Dadurch werden dann die
magnetischen Erscheinungen vollständig dargestellt, und für die meisten Personen
wird die Darstellung in dieser Form am angenehmsten sein. Allein
— was Ihnen vielleicht anfangs paradox scheinen wird — für denjenigen, der
versuchen will, das \textls[120]{Ganze der Erscheinungen} einer möglichst einfachen
Theorie unterwürfig zu machen, ist diese Darstellung nicht die zweckmässigste,
sondern eine andere wäre zu diesem Zweck von viel unmittelbarerer Brauchbarkeit.
Nemlich durch drei Karten, die die drei partiellen Intensitäten vor
Augen legten. Es sei \(m\) die ganze magnetische Kraft, \(i\) die Neigung, \(\delta\) die
Abweichung; dann werden die drei partiellen Kräfte:
\[
\begin{array}{@{}l@{\qquad}lllll@{}}
\xi=m\sin i & \text{in}&\text{verticaler}&\text{Richtung}&&\\
\eta=m\cos i\cos\delta\text{[\textsuperscript{*})]} &
\text{in}&\text{horizontaler}&\text{Richtung}&\text{nach}&\text{Norden}\\
\zeta=m\cos i\sin\delta\text{[\textsuperscript{*})]} &
\text{»}&\text{»}&\text{»}&\text{»}&\text{Westen.}
\end{array}
\]
\begingroup\renewcommand{\thefootnote}{}
\footnotetext{[\textsuperscript{*}) In der Handschrift steht \(d\) statt \(\delta\).]}\endgroup

\indent Wären die drei Karten für \(\xi,\eta,\zeta\) vorhanden, so wäre ich geneigt, einen
Versuch der oben angedeuteten Art zu machen; vielleicht entschlösse sich
Herr \textsc{Hansteen} dazu solche zu liefern, oder allenfalls auch nur Eine derselben.
Meine theoretische Untersuchung zeigt sogar, dass eine vollständige
Darstellung Einer partiellen Kraft an sich zureichend ist die andere a priori
abzuleiten. Selbst solche Karten erst zu entwerfen, werde ich mich nicht entschliessen,
da dazu eine längere innige critische Bekanntschaft mit den Quellen
erforderlich ist. Die Zurückführung auf eine kleine Anzahl von Polen z. B. 4,
halte ich übrigens für nicht naturgemäss; solche Pole sind nur Symptome in
den Erscheinungen, die keine scharfe Bedeutung haben, und wenn wir erst im
Besitz der allgemeinen alles auf einmal umfassenden Formel sind, ergeben
sich diese sogenannten Pole, wenn man sie wissen will, von selbst mit. Vielleicht
wird Ihnen, was ich sagen will, durch ein analoges Beispiel deutlicher.
Die Zeitgleichung bietet im Jahre mehrere Maxima und Minima dar, aber
man würde Unrecht haben, diesen eine ganz besondere Bedeutung beizulegen.

Mit einer anderen und wohl an sich nicht viel weniger wichtigen Seite
des Gegenstandes habe ich mich in den letzten Wochen viel und wie mir
deucht nicht ohne Erfolg beschäftigt, nemlich mit einem Mittel, die Intensität
\clearpage\SourceLeaf{80}\HistoricalHead{75}{R}{GAUSS AN SCHUMACHER.}\HistoricalFoot{}{10\textsuperscript{*}}
\noindent des Erdmagnetismus auf eine absolute Einheit zurückzuführen. Wenn ich
nicht irre, hat \textsc{Poisson}[\textsuperscript{*})]\begingroup\renewcommand{\thefootnote}{}
\footnotetext{[\textsuperscript{*}) \textsc{S. D. Poisson}, \textit{Solution d’un problème relatif au Magnétisme terrestre},
Additions à la Connaissance des tems pour 1828, Paris 1825, S. 322, Auszug in den Annales de Chimie
et de Physique \textbf{30}, 1825, S. 257, deutsch in \textsc{Baumgartners} Zeitschrift für Physik und Math.
\textbf{1}, Wien 1826, S. 117.]}
\endgroup\ zuerst ein Verfahren angegeben, und ich finde auch
in \textsc{Poggendorfs} Annalen[\textsuperscript{**}],\begingroup\renewcommand{\thefootnote}{}
\footnotetext{[\textsuperscript{**}) \textsc{Ludwig Moser} und \textsc{Peter Riess}, \textit{Über die Messung der Intensität
des tellurischen Magnetismus}, \textsc{Poggendorffs} Annalen der Physik und Chemie \textbf{18}, 1830, S. 226.]}
\endgroup\ einen Versuch, solches zur Anwendung zu bringen.
Allein ich finde dabei verschiedenes, was ich durchaus für unzulässig halten
muss, und halte mich überzeugt, dass durch solche Behandlung auch nicht
einmal ein grob genähertes Resultat erhalten werden kann. Ich habe mehrere
Reihen Versuche, aber unter andern Umständen, gemacht, deren schärfere
Berechnung, wie ich schon jetzt erkenne, eine ziemliche Annäherung geben
werden, deren Resultat aber weit von dem in \textsc{Poggendorfs} Annalen verschieden
ist (etwa \(\frac{1}{20}\) so gross\textsuperscript{***})).\begingroup\renewcommand{\thefootnote}{}
\footnotetext{\textsuperscript{***}) Der Unterschied erscheint noch viel greller, wenn man erwägt, daß die
Größe eigentlich aus ihrem Quadrate bestimmt wird, welches also dort 400mal zu
gross gefunden ist.}
\endgroup\ Allein ich bin auf ein anderes Verfahren gekommen,
welches ein viel reineres Resultat geben kann, und ich halte es für möglich,
selbst die Genauigkeit des Resultats, wenn man alle nöthigen Vorkehrungen
macht, so weit zu treiben, dass sie derjenigen, die durch vergleichende Beobachtungen
mit Einer Nadel [erzielt wird], an die Seite gestellt werden kann,
oder sie vielleicht noch überbietet. Schon jetzt geben die Versuche, die hauptsächlich
Freund \textsc{Weber} nach meinen Angaben gemacht hat, eine Genauigkeit,
worin wohl schwerlich mehr, als einige Procent Ungewissheit zurückbleiben;
man wird es aber viel weiter treiben können. Es ist gewiss in zwiefacher
Rücksicht sehr wichtig, dass wir hierin ins Klare kommen. Ist die Möglichkeit
erst da, wenn auch unter Anwendung von einigen Vorkehrungen, die
absolute Grösse des Erdmagnetismus zu bestimmen, so soll man sich dies an
einer Anzahl Örter über der ganzen Erde angelegen sein lassen; reisende
Beobachter führen invariable Nadeln bei sich, womit sie die Verhältnisse anderer
Örter unter sich bestimmen, und indem sie von Zeit zu Zeit solche Punkte
berühren, wo die absolute Intensität ausgemittelt ist, versichern sie sich der
bleibenden Invariabilität ihrer Nadeln, und führen ihre Resultate auf absolutes
Maass. Aber noch wichtiger ist es für künftige Jahr-
\clearpage\SourceLeaf{81}\HistoricalHead{76}{L}{PHYSIK. BRIEFWECHSEL.}
\noindent hunderte, in denen eben so bedeutende Änderungen in der absoluten Intensität
zu erwarten sind, wie wir lange bei der Declination und Neigung kennen.
Ich habe immer diese ungeheuren Änderungen, wie etwas höchst merkwürdiges
betrachtet. Ohne Zweifel ist die magnetische Erdkraft nicht das Resultat
von ein Paar grossen Magneten in der Nähe des Erdmittelpunkts,
die nach und nach viele Meilen weit sich von ihrem Platze bewegen, sondern
das Resultat aller in der Erde enthaltenen polarisirten Eisentheile, und zwar
mehr derjenigen, die der Oberfläche, als der, die dem Mittelpunkte näher liegen.
Allein was soll man von den ungeheuren Änderungen, die seit ein Paar Jahrhunderten
Statt gefunden haben, denken? Mir hat immer diese Erscheinung
eine besondere Gunst für die von \textsc{Cordier}[\textsuperscript{*})]\begingroup\renewcommand{\thefootnote}{}
\footnotetext{[\textsuperscript{*}) \textsc{L. Cordier}, \textit{Essai sur la Température de l’intérieur de la Terre},
Mémoires de l’Académie des Sciences \textbf{7}, Paris 1827, S. 473, insbesondere S. 538 und 540.]}
\endgroup\ besonders hervorgehobene Hypothese
zu erwecken geschienen, wonach die feste Erdrinde vergleichungsweise nur
dünn ist. Natürlich können dann nur in dieser die magnetischen Kräfte ihren
Sitz haben, und die allmähliche Verdickung dieser Rinde durch Erstarren vorher
flüssig gewesener Schichten erklärt dann die eintretende grosse Veränderung
in dem Erdmagnetismus auf das ungezwungenste, die sonst ein grosses Räthsel
bleibt. Auch der Umstand, dass die sogenannten magnetischen Hauptpole
der Erde in die kältesten Gegenden fallen, wo vermuthlich die Erdrinde am
dicksten ist, scheint darauf hinzudeuten.

Doch ich breche hier ab und bitte Sie, recht bald wieder mit einigen
Zeilen zu erfreuen
\begin{flushright}
Ihren ganzen eigenen\hspace{2em}\\[.5ex]
C. F. G[\textsc{auss}].
\end{flushright}
\noindent . . . . . . . . . . . . . . . . . . . . . . . . . . .\par
\indent Was in der ex post roth eingeklammerten Stelle[\textsuperscript{**})]\begingroup\renewcommand{\thefootnote}{}
\footnotetext{[\textsuperscript{**}) Gemeint sind die oben S. 75 in Klammern gesetzten Worte, »etwa
\(\frac{1}{20}\) so groß« und die zugehörige Fußnote.]}
\endgroup\ gesagt ist, bedarf
\textls[120]{vielleicht} einer Berichtigung; es ist in dem fraglichen Aufsatze nicht klar
ausgesprochen, was die Einheit eigentlich bedeutet, womit die magnetische
Intensität gemessen werden soll; ich vermuthe aber (ich habe jetzt nicht Zeit
dies gleich genauer zu untersuchen), dass, um die dortige Zahl mit der meinigen
\textls[120]{vergleichbar} zu machen, meine erst noch mit \(\surd\) (doppelte Fallzeit
\clearpage\SourceLeaf{82a}\HistoricalHead{77}{R}{GAUSS AN GERLING}
\noindent in Einer Secunde in der gewählten Einheit) multiplicirt werden muss; dieser
Factor ist nun nahe \(=99\), wonach dann mein Resultat gegen 5mahl so gross
ist, wie das dortige. Eine kleine Ungewissheit wird bei der Übersetzung immer
bleiben, da die Herren \textsc{Rieser} und \textsc{Moser}[\textsuperscript{*})]\begingroup\renewcommand{\thefootnote}{}
\footnotetext{[\textsuperscript{*}) So in der Handschrift; es soll heißen \textsc{Riess} und \textsc{Moser}, vergl. oben S. 75 die zweite Fußnote.]}
\endgroup\ ihre Nadeln nicht gewogen haben.
Eine freilich auch nur flüchtige Reduction meiner Messungen gibt, die Dichtigkeit
des Wassers \(=1\) gesetzt,
\begin{center}
\begin{tabular}{@{}l@{\quad}r@{\,}l@{}}
Intensität im horizontalen Sinn & 1,\textsuperscript{millim.}&72 aus den Schwingungsversuchen\\
&1,&59 aus den Versuchen nach der andern\\
&&\quad Methode.
\end{tabular}
\end{center}
\noindent\hspace*{\parindent}\textsc{Rieser} und \textsc{Moser} geben \(0{,}13017\), welches aber erst noch mit \(\surd\) Dichtigkeit
des Stahls multiplicirt werden muss; nehme ich diese \(=7{,}8163\), so kommt
\[
0{,}364.
\]
\clearpage\endgroup
